% ==============================================================================
% Apêndice: lista completa dos testes de robustez
% Incluir dentro do ambiente de apêndices do abntex2 (\begin{apendicesenv} ... \end{apendicesenv})
% As tabelas de resultados de cada teste estão em
% ../Modelagem/regressao/saidas/tabelas/robustez_R*.csv (e .docx)
% ==============================================================================

\chapter{Testes de robustez}\label{ap:robustez}

O \autoref{qua:robustez} lista todos os testes de robustez estimados. Cada um muda uma decisão em relação aos modelos principais; os testes R5, R10 e R11 mudam mais de um elemento, como indicado. Os quatro testes discutidos no texto principal são R5, R12, R14 e R16.

\begin{quadro}[htbp]
\caption{Testes de robustez}\label{qua:robustez}
\centering
\footnotesize
\begin{tabular}{|>{\raggedright\arraybackslash}p{1.0cm}|>{\raggedright\arraybackslash}p{8.0cm}|>{\raggedright\arraybackslash}p{5.0cm}|}
\hline
\textbf{Teste} & \textbf{O que muda em relação ao modelo principal} & \textbf{O que verifica} \\ \hline
R1 & Coalizões sem considerar fusões de partidos (só mudanças de nome) & Regra de padronização das siglas \\ \hline
R2 & Exclui município-anos com prefeito ou governador vindo de eleição suplementar & Casos atípicos de mandato \\ \hline
R3 & Variável dependente em logaritmo & Forma funcional (efeito em \%) \\ \hline
R4 & Variável dependente como participação na despesa total (\%) & Recomposição do gasto entre áreas \\ \hline
R5 & Especificação de \textcite{sakurai2009}: sem pré-eleitoral e sem pandemia, com erro agrupado por município & Comparação com a literatura \\ \hline
R6 & Inclui dummy de prefeito em segundo mandato & Controle de reeleição \\ \hline
R7 & PIB municipal no lugar do nacional (amostra 2013--2023) & Controle econômico local \\ \hline
R7a & Especificação principal só em 2013--2023 & Separa o efeito da amostra do efeito do PIB no R7 \\ \hline
R8 & Assistência Social somada à Previdência & Agrupamento original de \textcite{sakurai2009} \\ \hline
R9 & Transporte, Agricultura e Comunicações só com municípios que informam a função em todos os anos & Registro irregular das funções \\ \hline
R10 & Sem limpeza: mantém os município-anos atípicos e a receita tributária original & Regra de limpeza \\ \hline
R11 & Troca a exclusão dos atípicos pela winsorização das dependentes nos percentis 1 e 99 de cada ano & Sensibilidade a valores extremos \\ \hline
R12 & Modelo A com erro agrupado por município & Tipo de erro-padrão \\ \hline
R13 & Modelo B com erro de \textcite{driscoll1998} & Tipo de erro-padrão \\ \hline
R14 & Modelo A sem a dummy de pandemia & Efeito da dummy de 2020 \\ \hline
R15 & Sem as variáveis do Censo (jovens, idosos e urbanização) & Controles que variam pouco no tempo \\ \hline
R16 & Um efeito para cada eleição (2016 e 2024) e para cada ano pré-eleitoral & Se o efeito médio vem de uma só eleição \\ \hline
\end{tabular}
\fonte{Elaboração própria.}
\end{quadro}
